 It arrives through instruments, architectures, and digital systems that translate vibration into legible form. Forms with a specific history, a specific function, embedded and inherent. Sound is never neutral. Following Grimshaw and Garner’s idea of \textit{sonic virtuality}, I consider sound as a field of potential that comes into being through the conditions of its mediation \parencite{Grimshaw_Garner_2015}. Each microphone, algorithm, or interface performs a conversion of energy, turning sonic potentials into technical expressions. This conversion, a kind of \textit{transduction}, is always shaped by the histories and values embedded in tools themselves. As Don Ihde observes, technologies magnify and reduce experience in patterned ways, framing how the world can appear through them \parencite{Ihde_1990}. Ruha Benjamin’s discussion of technological “default settings” extends this point to the social domain, showing how design embeds racialized and cultural assumptions that become invisible through habitual use \parencite{Benjamin_2019a}.

In high school, as a young musician, it was explained to me by my assistant band director that my trombone has a specific temperament. My Yamaha produced a slightly sharper tone in the typical "third position" area of the slide (when tuned to a standard pitch), different from that of a Bach or Shire model. Each maker, he said, designed, built, and tuned the instrument according to a distinct philosophy of sound. That detail, which is just mere millimeters of brass, taught me that instruments carry with them the intentions and sonic ideals of their makers. Now, I recognize that moment as a lesson in epistemology: every device, from a trombone to a software synthesizer, expresses a particular worldview about what counts as correct, beautiful, or possible.

This dissertation begins from such recognition. To modify a tool is to engage in an epistemic act: to re-tune inherited assumptions and redirect the flow of knowledge. Through the framework of \textit{Speculocultural Technopoiesis}, I treat modification as a process of transduction—where speculative cultural imagination takes material form—and refraction—where altered instruments bend perception toward new sonic and cultural relationships. The research situates these acts within the intertwined legacies of \textit{sonic virtuality} \parencite{Grimshaw_Garner_2015}, \textit{sonic fiction} \parencite{Eshun_1999}, and culturally grounded practices of making, positioning the sonic as both medium and method for re-imagining technological knowledge.



